\documentclass[10pt,a4paper]{amsart}
\usepackage[margin=19mm]{geometry}
\usepackage{amsmath,amssymb,mathtools}
\usepackage[scr=rsfs,cal=euler]{mathalfa}
\usepackage{xfrac}
\usepackage[unicode,hidelinks,bookmarks=false]{hyperref}
\newcommand{\ie}{i.e.\ }
\newcommand{\cf}{cf.\ }
\newcommand{\resp}{respectively\ }
\newcommand{\Subsection}{\subsection}
\providecommand{\nobreakdash}{\nobreak\hbox{-}}
\setlength{\emergencystretch}{2em}
\multlinegap0pt
\usepackage{mathtools}
\usepackage{amssymb}
\newtheorem{prop}{}
\newtheorem{thm}{}
\title{The universal logic of repeated experiments}
\author{Sergio Daniel Grillo}
\date{Source: arXiv:2601.00118v2 (CC BY 4.0)}
\begin{document}
\maketitle

\noindent\emph{Source and licence.} The universal logic of repeated experiments. Version arXiv:2601.00118v2 (CC BY 4.0), distributed by arXiv under the Creative Commons Attribution 4.0 International licence, http://creativecommons.org/licenses/by/4.0/, as recorded in the arXiv licence metadata for this exact version and shown on the abstract page https://arxiv.org/abs/2601.00118v2. Full licence notice: \texttt{licenses/arxiv-2601.00118-CC-BY-4.0.txt}.\par\smallskip

\noindent\textit{Editorial scope.} Counted unit: the article's thm tia (source0.tex lines 2326--2332). The article's complete original proof of it is delivered with it (source0.tex lines 2334--2464, one \texttt{proof} environment of 131 lines). No mathematics is written, repaired or re-derived: the statement and the proof are the article's own lines, in the article's own notation, and no retained line is substituted or re-flowed.  Retained uncounted as the source-local context the proof needs: lines 1972--1974; lines 1991--1993; lines 2050--2077; lines 2314--2316.  Nothing was omitted from any retained span, so the package carries no omission record.  No retained line contains a citation, so the wrapper carries no bibliography and no bibliography entry.  No retained line contains a figure environment, an image-inclusion command or drawing code, so no image is delivered and the package declares no asset dependency.  Editorial formatting only, in addition: an AMS \texttt{amsart} preamble that restates the article's own declarations and macros used by the retained lines, namely the environments \texttt{prop}, \texttt{thm} on separate counters, together with the standard packages those lines need.  The retained environments print in this document as Proposition 1, Theorem 1 in order of appearance, whereas the article prints the same items with the numbering of its own full document.  The complete native source of the article as distributed in the arXiv e-print for this exact version is bundled unchanged as \texttt{sources/191972/source0.tex}.
\begin{equation}
\varphi:A\in\mathsf{\mathsf{F}_{\kappa}}\mapsto\left[A\right]\in\mathsf{D}\left(\mathsf{F}_{\kappa}\right).\label{fial}
\end{equation}

\begin{equation}
{\textstyle \left(\bigvee_{i\in I}A^{i}\right)^{*}={\textstyle \bigvee_{f\in F}}\mathbb{A}_{f,I}},\quad\mathbb{A}_{f,I}\left(\alpha\right)\coloneqq{\textstyle \bigcap_{f\left(i\right)=\alpha}}\left(A^{i}\left(\alpha\right)\right)^{c},\label{fc}
\end{equation}

Related to each $\alpha$, we have a natural embedding
\[
\mathfrak{i}_{\alpha}:a\in\mathsf{E}_{\alpha}\longmapsto\mathfrak{i}_{\alpha}\left(a\right)\in\mathsf{F}_{\kappa}
\]
with
\begin{equation}
\left(\mathfrak{i}_{\alpha}\left(a\right)\right)\left(\beta\right)=\begin{cases}
a, & \beta=\alpha,\\
\mathbf{1}, & \beta\neq\alpha,
\end{cases}\label{ia}
\end{equation}
which, \textit{via} the embedding $\varphi:\mathsf{F}_{\kappa}\rightarrow{\textstyle \bigotimes_{\beta\in\kappa}}\mathsf{E}_{\beta}$
(see \eqref{fial}), can be seen as a function from $\mathsf{E}_{\alpha}$
into ${\textstyle \bigotimes_{\beta\in\kappa}}\mathsf{E}_{\beta}$.
Below, by a \textit{logic embedding} we mean an injective homomorphism
of logics (i.e. of orthocomplemented complete lattices).
\begin{prop}
Each function $\mathfrak{i}_{\alpha}:\mathsf{E}_{\alpha}\rightarrow{\textstyle \bigotimes_{\beta\in\kappa}}\mathsf{E}_{\beta}$,
given by Eq. \eqref{ia}, is a logic embedding. In particular,
\[
\mathfrak{i}_{\alpha}\left(\bigcup S\right)=\coprod\mathfrak{i}_{\alpha}\left(S\right),\quad\mathfrak{i}_{\alpha}\left(\bigcap S\right)=\bigwedge\mathfrak{i}_{\alpha}\left(S\right),
\]
for all $S\subseteq\mathsf{E}_{\alpha}$, and 
\[
\mathfrak{i}_{\alpha}\left(\mathbf{0}\right)=\hat{\mathbf{0}},\quad\mathfrak{i}_{\alpha}\left(\mathbf{1}\right)=\hat{\mathbf{1}},\quad\mathfrak{i}_{\alpha}\left(a^{c}\right)=\left(\mathfrak{i}_{\alpha}\left(a\right)\right)^{*},
\]
for all $a\in\mathsf{E}_{\alpha}$.
\end{prop}

\begin{equation}
{\textstyle \bigwedge_{i\in I}\bigvee_{\alpha\in\kappa}}e_{\alpha}\left(B^{i}\left(\alpha\right)\right)={\textstyle \bigvee_{f\in F}\bigwedge_{\alpha\in\kappa}}e_{\alpha}\left({\textstyle \bigcap_{f\left(i\right)=\alpha}}B^{i}\left(\alpha\right)\right).\label{edb}
\end{equation}

\begin{thm}
There exists a unique homomorphism of logics $t:{\textstyle \bigotimes_{\alpha\in\kappa}}\mathsf{E}_{\alpha}\rightarrow\mathsf{T}$
such that 
\begin{equation}
t\left(\mathfrak{i}_{\alpha}\left(a\right)\right)=e_{\alpha}\left(a\right),\quad\forall a\in\mathsf{E}_{\alpha},\quad\forall\alpha\in\kappa.\label{tia}
\end{equation}
\end{thm}

\begin{proof}
Consider the function $\tilde{t}:\mathsf{D}\left(\mathsf{F}_{\kappa}\right)\rightarrow\mathsf{T}$
such that 
\begin{equation}
\tilde{t}\left({\textstyle \bigvee_{i\in I}}A^{i}\right)={\textstyle \bigvee_{i\in I}}{\textstyle \bigwedge_{\alpha\in\kappa}}e_{\alpha}\left(A^{i}\left(\alpha\right)\right).\label{dt}
\end{equation}
Let us see $\tilde{t}$ is well-defined. If ${\textstyle \bigvee_{i\in I}}A^{i}={\textstyle \bigvee_{j\in J}}B^{j}$,
this means that for all $i\in I$ there exists $j\left(i\right)\in J$
such that $A^{i}\subseteq B^{j\left(i\right)}$, i.e. 
\[
A^{i}\left(\alpha\right)\subseteq B^{j\left(i\right)}\left(\alpha\right),\quad\forall\alpha\in\kappa,
\]
and vice versa (changing $A$'s by $B$'s). So, 
\[
\begin{array}{lll}
\tilde{t}\left({\textstyle \bigvee_{i\in I}}A^{i}\right) & = & {\textstyle \bigvee_{i\in I}}{\textstyle \bigwedge_{\alpha\in\kappa}}e_{\alpha}\left(A^{i}\left(\alpha\right)\right)\leq{\textstyle \bigvee_{i\in I}}{\textstyle \bigwedge_{\alpha\in\kappa}}e_{\alpha}\left(B^{j\left(i\right)}\left(\alpha\right)\right)\\
\\
 & \leq & {\textstyle \bigvee_{j\in J}}{\textstyle \bigwedge_{\alpha\in\kappa}}e_{\alpha}\left(B^{j}\left(\alpha\right)\right)=\tilde{t}\left({\textstyle \bigvee_{j\in J}}B^{j}\right),
\end{array}
\]
and the other inequality follows by changing $A$'s by $B$'s. 

Note that, if $A\in\mathsf{F}_{\kappa}$, then we can write $A=\bigvee_{i\in I}A^{i}$
with $I=\left\{ s\right\} $ and $A^{s}=A$. Thus, according to \eqref{dt},
\begin{equation}
\tilde{t}\left(A\right)={\textstyle \bigwedge_{\alpha\in\kappa}}e_{\alpha}\left(A\left(\alpha\right)\right).\label{dt1}
\end{equation}
In particular, given $a\in\mathsf{E}_{\beta}$,
\[
\tilde{t}\left(\mathfrak{i}_{\beta}\left(a\right)\right)={\textstyle \bigwedge_{\alpha\in\kappa}}e_{\alpha}\left(\left(\mathfrak{i}_{\beta}\left(a\right)\right)\left(\alpha\right)\right).
\]
Since $\left(\mathfrak{i}_{\beta}\left(a\right)\right)\left(\alpha\right)=\mathbf{1}$
when $\alpha\neq\beta$ (see \eqref{ia}), then 
\begin{equation}
\tilde{t}\left(\mathfrak{i}_{\beta}\left(a\right)\right)=e_{\beta}\left(\left(\mathfrak{i}_{\beta}\left(a\right)\right)\left(\beta\right)\right)=e_{\beta}\left(a\right),\quad\forall a\in\mathsf{E}_{\beta},\quad\forall\beta\in\kappa,\label{tia2}
\end{equation}
because $\left(\mathfrak{i}_{\beta}\left(a\right)\right)\left(\beta\right)=a$
(see \eqref{ia} again).

From \eqref{dt} and \eqref{dt1}, we have that
\begin{equation}
\tilde{t}\left({\textstyle \bigvee_{i\in I}}A^{i}\right)={\textstyle \bigvee_{i\in I}}\tilde{t}\left(A^{i}\right),\quad\forall\;{\textstyle \bigvee_{i\in I}}A^{i}\in\mathsf{D}\left(\mathsf{F}_{\kappa}\right),\label{tas}
\end{equation}
and given a subset 
\[
\left\{ \mathbb{A}_{j}:j\in J\right\} \subseteq\mathsf{D}\left(\mathsf{F}_{\kappa}\right),
\]
with $\mathbb{A}_{j}={\textstyle \bigvee_{i\in I_{j}}}A_{j}^{i}$,
\begin{equation}
\begin{array}{lll}
\tilde{t}\left({\textstyle \bigvee_{j\in J}}\mathbb{A}_{j}\right) & = & \tilde{t}\left({\textstyle \bigvee_{j\in J}}{\textstyle \bigvee_{i\in I_{j}}}A_{j}^{i}\right)={\textstyle \bigvee_{j\in J}}{\textstyle \bigvee_{i\in I_{j}}}\tilde{t}\left(A_{j}^{i}\right)\\
\\
 & = & {\textstyle \bigvee_{j\in J}}\tilde{t}\left({\textstyle \bigvee_{i\in I_{j}}}A_{j}^{i}\right)={\textstyle \bigvee_{j\in J}}\tilde{t}\left(\mathbb{A}_{j}\right).
\end{array}\label{ttu}
\end{equation}
Also, from \eqref{fc}, \eqref{dt1} and \eqref{tas},
\[
\tilde{t}\left(\left({\textstyle \bigvee_{i\in I}}A^{i}\right)^{*}\right)=\tilde{t}\left({\textstyle \bigvee_{f\in F}}\mathbb{A}_{f,I}\right)={\textstyle \bigvee_{f\in F}}\tilde{t}\left(\mathbb{A}_{f,I}\right)
\]
and
\[
\tilde{t}\left(\mathbb{A}_{f,I}\right)={\textstyle \bigwedge_{\alpha\in\kappa}}e_{\alpha}\left(\mathbb{A}_{f,I}\left(\alpha\right)\right)={\textstyle \bigwedge_{\alpha\in\kappa}}e_{\alpha}\left({\textstyle \bigcap_{f\left(i\right)=\alpha}}\left(A^{i}\left(\alpha\right)\right)^{c}\right).
\]
On the other hand, the meet-join distributivity property for $\left(e,\mathsf{T}\right)$
ensures that\footnote{See Eq. \eqref{edb} for $B^{i}\left(\alpha\right)=\left(A^{i}\left(\alpha\right)\right)^{c}$.}
\[
{\textstyle \bigvee_{f\in F}}{\textstyle \bigwedge_{\alpha\in\kappa}}e_{\alpha}\left({\textstyle \bigcap_{f\left(i\right)=\alpha}}\left(A^{i}\left(\alpha\right)\right)^{c}\right)={\textstyle \bigwedge_{i\in I}\bigvee_{\alpha\in\kappa}}e_{\alpha}\left(\left(A^{i}\left(\alpha\right)\right)^{c}\right)=\left({\textstyle \bigvee_{i\in I}\bigwedge_{\alpha\in\kappa}}e_{\alpha}\left(A^{i}\left(\alpha\right)\right)\right)^{\bullet},
\]
so
\begin{equation}
\tilde{t}\left(\left({\textstyle \bigvee_{i\in I}}A^{i}\right)^{*}\right)=\left({\textstyle \bigvee_{i\in I}\bigwedge_{\alpha\in\kappa}}e_{\alpha}\left(A^{i}\left(\alpha\right)\right)\right)^{\bullet}=\left(\tilde{t}\left({\textstyle \bigvee_{i\in I}}A^{i}\right)\right)^{\bullet}.\label{tst}
\end{equation}

Now, define $t:{\textstyle \bigotimes_{\alpha\in\kappa}}\mathsf{E}_{\alpha}\rightarrow\mathsf{T}$
as the restriction of $\tilde{t}$. Since $\mathsf{F}_{\kappa}\subseteq{\textstyle \bigotimes_{\alpha\in\kappa}}\mathsf{E}_{\alpha}$,
it is clear from \eqref{tia2} that $t$ satisfies \eqref{tia}. Also,
given ${\textstyle \bigvee_{i\in I}}A^{i}\in{\textstyle \bigotimes_{\alpha\in\kappa}}\mathsf{E}_{\alpha}$,
since $\left({\textstyle \bigvee_{i\in I}}A^{i}\right)^{*}\in{\textstyle \bigotimes_{\alpha\in\kappa}}\mathsf{E}_{\alpha}$
too, then \eqref{tst} implies that
\[
t\left(\left({\textstyle \bigvee_{i\in I}}A^{i}\right)^{*}\right)=\left(t\left({\textstyle \bigvee_{i\in I}}A^{i}\right)\right)^{\bullet},
\]
i.e. $t$ respect the orthocomplementation. Finally, given an arbitrary
subset 
\[
\left\{ \mathbb{A}_{j}:j\in J\right\} \subseteq{\textstyle \bigotimes_{\alpha\in\kappa}}\mathsf{E}_{\alpha},
\]
from \eqref{ttu} and \eqref{tst} we have that
\[
\begin{array}{lll}
t\left({\textstyle \coprod_{j\in J}}\mathbb{A}_{j}\right) & = & t\left(\left({\textstyle \bigvee_{j\in J}}\mathbb{A}_{j}\right)^{**}\right)=\tilde{t}\left(\left({\textstyle \bigvee_{j\in J}}\mathbb{A}_{j}\right)^{**}\right)=\left(\tilde{t}\left({\textstyle \bigvee_{j\in J}}\mathbb{A}_{j}\right)\right)^{\bullet\bullet}\\
\\
 & = & \tilde{t}\left({\textstyle \bigvee_{j\in J}}\mathbb{A}_{j}\right)={\textstyle \bigvee_{j\in J}}\tilde{t}\left(\mathbb{A}_{j}\right)={\textstyle \bigvee_{j\in J}}t\left(\mathbb{A}_{j}\right).
\end{array}
\]
The fact that $t$ respects the meets follows from the De Morgan laws.
Summing up, the function $t:{\textstyle \bigotimes_{\alpha\in\kappa}}\mathsf{E}_{\alpha}\rightarrow\mathsf{T}$
such that
\begin{equation}
t\left({\textstyle \bigvee_{i\in I}}A^{i}\right)={\textstyle \bigvee_{i\in I}}{\textstyle \bigwedge_{\alpha\in\kappa}}e_{\alpha}\left(A^{i}\left(\alpha\right)\right)\label{dt2}
\end{equation}
 is a homomorphism of logics that satisfies \eqref{tia}. It rests
to see that it is the only one.

\bigskip{}

Note first that any element $A\in\mathsf{F}_{\kappa}$ satisfies
\begin{equation}
A={\textstyle \bigwedge_{\alpha\in\kappa}}\mathfrak{i}_{\alpha}\left(A\left(\alpha\right)\right).\label{aA}
\end{equation}
In fact, ${\textstyle \bigwedge_{\alpha\in\kappa}}\mathfrak{i}_{\alpha}\left(A\left(\alpha\right)\right)={\textstyle \bigcap_{\alpha\in\kappa}}\mathfrak{i}_{\alpha}\left(A\left(\alpha\right)\right)\in\mathsf{F}_{\kappa}$
and
\[
\left({\textstyle \bigcap_{\alpha\in\kappa}}\mathfrak{i}_{\alpha}\left(A\left(\alpha\right)\right)\right)\left(\beta\right)={\textstyle \bigcap_{\alpha\in\kappa}}\left(\mathfrak{i}_{\alpha}\left(A\left(\alpha\right)\right)\left(\beta\right)\right)=\left(\mathfrak{i}_{\beta}\left(A\left(\beta\right)\right)\left(\beta\right)\right)=A\left(\beta\right).
\]
Suppose that $t':{\textstyle \bigotimes_{\alpha\in\kappa}}\mathsf{E}_{\alpha}\rightarrow\mathsf{T}$
is a homomorphism of logics satisfying \eqref{tia}. If ${\textstyle \bigvee_{i\in I}}A^{i}\in{\textstyle \bigotimes_{\alpha\in\kappa}}\mathsf{E}_{\alpha}$,
then
\[
{\textstyle \bigvee_{i\in I}}A^{i}={\textstyle \coprod_{i\in I}}A^{i}={\textstyle \coprod_{i\in I}}{\textstyle \bigwedge_{\alpha\in\kappa}}\mathfrak{i}_{\alpha}\left(A^{i}\left(\alpha\right)\right)
\]
and (recall \eqref{dt})
\[
\begin{array}{lll}
t'\left({\textstyle \bigvee_{i\in I}}A^{i}\right) & = & t'\left({\textstyle \coprod_{i\in I}}A^{i}\right)=t'\left({\textstyle \coprod_{i\in I}}{\textstyle \bigwedge_{\alpha\in\kappa}}\mathfrak{i}_{\alpha}\left(A^{i}\left(\alpha\right)\right)\right)={\textstyle \bigvee_{i\in I}}{\textstyle \bigwedge_{\alpha\in\kappa}}t'\left(\mathfrak{i}_{\alpha}\left(A^{i}\left(\alpha\right)\right)\right)\\
\\
 & = & {\textstyle \bigvee_{i\in I}}{\textstyle \bigwedge_{\alpha\in\kappa}}e_{\alpha}\left(A^{i}\left(\alpha\right)\right)=t\left({\textstyle \bigvee_{i\in I}}A^{i}\right),
\end{array}
\]
what ensures that $t'=t$, as wanted. 
\end{proof}

\end{document}
